% ═══════════════════════════════════════════════════════════════════
% Mental Health Text Classification Using Fine-Tuned RoBERTa-Large
% with Focal Loss, Triple Pooling, and Temperature Scaling
% ═══════════════════════════════════════════════════════════════════
\documentclass[conference, 10pt]{IEEEtran}

\usepackage[utf8]{inputenc}
\usepackage[T1]{fontenc}
\usepackage{amsmath, amssymb}
\usepackage{graphicx}
\usepackage{booktabs}
\usepackage{multirow}
\usepackage{hyperref}
\usepackage{xcolor}
\usepackage{float}
\usepackage{enumitem}
\usepackage{algorithm}
\usepackage{algpseudocode}
\usepackage{caption}
\usepackage{subcaption}
\usepackage{array}

\hypersetup{
    colorlinks=true,
    linkcolor=blue,
    citecolor=blue,
    urlcolor=blue,
}

\title{Mental Health Text Classification Using Fine-Tuned RoBERTa-Large\\with Focal Loss, Triple Pooling, and Temperature Scaling}

% ── Replace with your actual author details
\author{
    \IEEEauthorblockN{Your Name Here}
    \IEEEauthorblockA{
        Department of Computer Science\\
        Your University Name\\
        City, Country\\
        email@example.com
    }
}

\begin{document}

\maketitle

% ═══════════════════════════════════════════════════════════════
% ABSTRACT
% ═══════════════════════════════════════════════════════════════
\begin{abstract}
Mental health disorders represent a growing global concern, and early detection through natural language analysis holds significant promise for scalable screening. This paper presents a fine-tuned RoBERTa-large transformer model for multi-class mental health text classification, capable of categorizing text into seven classes: Normal, Depression, Suicidal, Anxiety, Bipolar, Stress, and Personality Disorder. We propose a classification pipeline that incorporates \textbf{(i)}~partial fine-tuning with layer-wise learning rate decay, \textbf{(ii)}~a triple pooling strategy combining CLS, mean, and max pooling for richer representations, \textbf{(iii)}~focal loss with label smoothing to address class imbalance, and \textbf{(iv)}~post-training temperature scaling for improved probability calibration. Trained on a dataset of 53,043 samples with back-translation augmentation, the model achieves a test accuracy of \textbf{91.67\%} and a macro F1-score of \textbf{90.74\%} across all seven classes. We also present an interactive Streamlit-based web application for real-time inference. Experimental results demonstrate that the combination of partial fine-tuning, triple pooling, and focal loss significantly outperforms standard fine-tuning baselines on this imbalanced mental health dataset.
\end{abstract}

\begin{IEEEkeywords}
Natural Language Processing, Mental Health, Text Classification, RoBERTa, Transformer, Focal Loss, Transfer Learning, Deep Learning
\end{IEEEkeywords}

% ═══════════════════════════════════════════════════════════════
% I. INTRODUCTION
% ═══════════════════════════════════════════════════════════════
\section{Introduction}

Mental health disorders affect approximately one in eight people worldwide, according to the World Health Organization (WHO)~\cite{who2022}. Depression, anxiety, bipolar disorder, and suicidal ideation remain among the most prevalent conditions, yet access to timely diagnosis and treatment continues to be a significant challenge, particularly in resource-constrained settings. The proliferation of digital communication platforms—including social media, online forums, and messaging services—has created opportunities for automated screening systems that analyze textual expressions for mental health indicators~\cite{calvo2017natural}.

Natural Language Processing (NLP) has emerged as a powerful tool for mental health analysis. Early approaches relied on hand-crafted lexical features and traditional machine learning classifiers such as Support Vector Machines and Random Forests~\cite{coppersmith2015clpsych}. The advent of pre-trained transformer models—BERT~\cite{devlin2019bert}, RoBERTa~\cite{liu2019roberta}, and their variants—has dramatically improved performance on text classification tasks by capturing deep contextual relationships within language.

However, deploying transformer models for mental health classification presents unique challenges:
\begin{itemize}[noitemsep]
    \item \textbf{Class imbalance}: Mental health datasets are inherently imbalanced, with overrepresentation of depression and underrepresentation of conditions like personality disorder.
    \item \textbf{Overfitting risk}: Large models (355M+ parameters) are prone to overfitting on moderately-sized datasets.
    \item \textbf{Overconfidence}: Standard softmax probabilities are often poorly calibrated, producing overconfident predictions.
    \item \textbf{Subtle distinctions}: Overlapping symptoms across mental health conditions make multi-class discrimination challenging.
\end{itemize}

To address these challenges, this paper proposes a comprehensive classification pipeline that fine-tunes RoBERTa-large with several key innovations:

\begin{enumerate}[noitemsep]
    \item \textbf{Partial fine-tuning} with layer-wise learning rate decay ($\gamma = 0.9$) to prevent catastrophic forgetting of pre-trained knowledge while adapting to the mental health domain.
    \item \textbf{Triple pooling} (CLS + mean + max) to create a 3,072-dimensional representation that captures complementary aspects of the input text.
    \item \textbf{Focal loss} with class weighting and label smoothing to handle severe class imbalance and prevent overconfident predictions.
    \item \textbf{Temperature scaling} as a post-training calibration step to improve probability reliability.
    \item \textbf{Back-translation augmentation} for minority classes to enrich training data diversity.
\end{enumerate}

The resulting model achieves 91.67\% accuracy and 90.74\% macro F1-score on a held-out test set of 53,043 samples spanning seven mental health categories. An interactive web application built with Streamlit provides real-time inference, batch analysis, and model interpretability features.

The remainder of this paper is organized as follows: Section~II reviews related work. Section~III details the proposed methodology. Section~IV describes the experimental setup. Section~V presents and analyzes results. Section~VI discusses the deployed application. Section~VII covers limitations and ethical considerations, and Section~VIII concludes the paper.


% ═══════════════════════════════════════════════════════════════
% II. RELATED WORK
% ═══════════════════════════════════════════════════════════════
\section{Related Work}

\subsection{NLP for Mental Health Detection}

The intersection of NLP and mental health has attracted significant research interest. Early work by Coppersmith et al.~\cite{coppersmith2015clpsych} demonstrated that linguistic signals in social media posts correlate with mental health conditions. De Choudhury et al.~\cite{dechoudhury2013predicting} used behavioral attributes from Twitter to predict depression onset. These initial studies relied on bag-of-words features, TF-IDF, and LIWC (Linguistic Inquiry and Word Count) features combined with traditional classifiers.

The CLPsych shared tasks~\cite{zirikly2019clpsych} established benchmarks for assessing suicide risk from Reddit posts, spurring development of more sophisticated NLP approaches. Matero et al.~\cite{matero2019suicide} achieved strong performance using contextual embeddings from BERT, demonstrating the superiority of transformer-based approaches.

\subsection{Pre-trained Language Models}

BERT~\cite{devlin2019bert} introduced bidirectional pre-training for language representations, achieving state-of-the-art performance across NLP benchmarks. RoBERTa~\cite{liu2019roberta} improved upon BERT through optimized pre-training—dynamic masking, larger batch sizes, and more training data—resulting in superior downstream task performance. For mental health applications, Ji et al.~\cite{ji2022mentalbert} introduced MentalBERT, pre-trained on mental health forum data, showing domain-specific pre-training can further improve classification.

\subsection{Handling Class Imbalance}

Class imbalance is a persistent challenge in mental health classification. Traditional oversampling (SMOTE) and undersampling methods have limited effectiveness for high-dimensional text representations. Lin et al.~\cite{lin2017focal} proposed focal loss for object detection, which has since been adapted for text classification to down-weight easy examples and focus learning on hard cases. Combined with class-frequency-based weighting, focal loss provides a principled approach to imbalanced classification~\cite{cui2019class}.

\subsection{Pooling Strategies}

Standard transformer-based classifiers typically use the CLS token representation. However, recent work~\cite{reimers2019sentencebert} demonstrated that mean pooling over all token representations often outperforms CLS-only approaches for sentence-level tasks. Our work extends this by concatenating CLS, mean, and max pooling to form a richer representation, an approach inspired by successful applications in natural language inference~\cite{conneau2017supervised}.


% ═══════════════════════════════════════════════════════════════
% III. METHODOLOGY
% ═══════════════════════════════════════════════════════════════
\section{Methodology}

\subsection{Problem Formulation}

Given a text input $\mathbf{x} = (x_1, x_2, \ldots, x_n)$, the goal is to predict the mental health category $y \in \{1, 2, \ldots, 7\}$ corresponding to: Normal, Depression, Suicidal, Anxiety, Bipolar, Stress, and Personality Disorder. The model outputs a probability distribution $\hat{p}(y|\mathbf{x}) \in \mathbb{R}^7$ obtained via temperature-scaled softmax.

\subsection{Architecture Overview}

The proposed architecture consists of four major components, illustrated in Figure~\ref{fig:architecture}:

\begin{figure}[t]
    \centering
    \includegraphics[width=\columnwidth]{architecture.jpg}
    \caption{End-to-end model architecture: Input text is tokenized, processed through a partially fine-tuned RoBERTa-large encoder, aggregated via triple pooling, and classified through a three-layer head with temperature-scaled softmax output.}
    \label{fig:architecture}
\end{figure}

\subsubsection{Tokenization}
Input text is tokenized using RoBERTa's Byte-Pair Encoding (BPE) tokenizer with a maximum sequence length of 256 tokens. Sequences are padded or truncated to produce fixed-size tensors of input IDs and attention masks.

\subsubsection{Encoder — RoBERTa-large}
The encoder is a 24-layer RoBERTa-large transformer with a hidden dimension of 1,024 and approximately 355 million parameters. To balance knowledge retention and task adaptation:

\begin{itemize}[noitemsep]
    \item \textbf{Layers 0--13 (14 layers)}: Frozen — preserves general linguistic knowledge.
    \item \textbf{Layers 14--23 (10 layers)}: Trainable — adapts to mental health domain.
    \item \textbf{Embedding layer}: Frozen — tokenization-level representations are kept intact.
\end{itemize}

This partial fine-tuning strategy reduces trainable parameters to approximately 134.6M (38\% of total), significantly mitigating overfitting risk on a ~52K sample dataset.

\subsubsection{Triple Pooling}
Instead of using only the CLS token, we concatenate three complementary pooling operations on the final hidden states $\mathbf{H} \in \mathbb{R}^{n \times 1024}$:

\begin{align}
    \mathbf{v}_{\text{CLS}} &= \mathbf{H}[0] \in \mathbb{R}^{1024} \label{eq:cls} \\
    \mathbf{v}_{\text{mean}} &= \frac{\sum_{i=1}^{n} m_i \cdot \mathbf{H}[i]}{\sum_{i=1}^{n} m_i} \in \mathbb{R}^{1024} \label{eq:mean} \\
    \mathbf{v}_{\text{max}} &= \max_{i: m_i = 1} \mathbf{H}[i] \in \mathbb{R}^{1024} \label{eq:max} \\
    \mathbf{v}_{\text{pool}} &= [\mathbf{v}_{\text{CLS}}; \mathbf{v}_{\text{mean}}; \mathbf{v}_{\text{max}}] \in \mathbb{R}^{3072} \label{eq:concat}
\end{align}

where $m_i$ is the attention mask for token $i$. This 3,072-dimensional vector captures the overall sentence representation (CLS), average semantics (mean), and the most salient features (max).

\subsubsection{Classification Head}

The pooled representation is passed through a three-layer classification head:

\begin{equation}
\begin{aligned}
    \mathbf{z}_1 &= \text{Dropout}_{0.30}(\text{GELU}(\mathbf{W}_1 \cdot \text{LayerNorm}(\mathbf{v}_{\text{pool}}) + \mathbf{b}_1)) \\
    \mathbf{z}_2 &= \text{Dropout}_{0.15}(\text{GELU}(\mathbf{W}_2 \cdot \mathbf{z}_1 + \mathbf{b}_2)) \\
    \hat{\mathbf{y}} &= \mathbf{W}_3 \cdot \mathbf{z}_2 + \mathbf{b}_3
\end{aligned}
\end{equation}

where $\mathbf{W}_1 \in \mathbb{R}^{512 \times 3072}$, $\mathbf{W}_2 \in \mathbb{R}^{128 \times 512}$, $\mathbf{W}_3 \in \mathbb{R}^{7 \times 128}$. GELU activation and progressive dropout (0.30 → 0.15) provide non-linearity and regularization.

\subsection{Loss Function — Focal Loss with Label Smoothing}

To address class imbalance, we employ focal loss~\cite{lin2017focal} with class-frequency weighting and label smoothing:

\begin{equation}
    \mathcal{L}_{\text{focal}} = -\frac{1}{N} \sum_{i=1}^{N} \alpha_{y_i} \cdot (1 - p_{t,i})^{\gamma} \cdot \text{CE}_{\text{smooth}}(\hat{\mathbf{y}}_i, y_i)
    \label{eq:focal}
\end{equation}

where:
\begin{itemize}[noitemsep]
    \item $p_{t,i}$ is the predicted probability for the true class
    \item $\gamma = 2.0$ is the focusing parameter that down-weights easy examples
    \item $\alpha_{y_i} = \sqrt{1/f_{y_i}}$ is the class weight based on inverse square root frequency, normalized so that $\sum \alpha_c = C$
    \item $\text{CE}_{\text{smooth}}$ applies label smoothing with $\epsilon = 0.05$
\end{itemize}

The smoothed target distribution is:
\begin{equation}
    q_c = \begin{cases}
        1 - \epsilon & \text{if } c = y \\
        \frac{\epsilon}{C - 1} & \text{otherwise}
    \end{cases}
\end{equation}

\subsection{Optimization Strategy}

\subsubsection{Layer-wise Learning Rate Decay}
Different learning rates are assigned to each trainable transformer layer. The top layer (layer 23) receives the base learning rate $\eta_{\text{bert}} = 2 \times 10^{-5}$. Each layer below is scaled by a decay factor $\gamma = 0.9$:

\begin{equation}
    \eta_{\ell} = \eta_{\text{bert}} \cdot \gamma^{(L - 1 - \ell)}
    \label{eq:lrdecay}
\end{equation}

where $\ell$ is the layer index among the $L = 10$ trainable layers. The classification head uses a separate, higher learning rate $\eta_{\text{head}} = 1 \times 10^{-4}$.

\subsubsection{Training Procedure}
\begin{itemize}[noitemsep]
    \item \textbf{Optimizer}: AdamW with weight decay $\lambda = 0.02$
    \item \textbf{Scheduler}: Cosine annealing with 6\% linear warmup
    \item \textbf{Gradient accumulation}: 4 steps (effective batch size = 64)
    \item \textbf{Mixed precision}: FP16 for memory efficiency on GPU
    \item \textbf{Gradient clipping}: Max norm = 1.0
    \item \textbf{Early stopping}: Patience of 5 epochs on validation macro F1
    \item \textbf{Maximum epochs}: 20
\end{itemize}

\subsection{Post-training Temperature Scaling}

After training, we apply temperature scaling~\cite{guo2017calibration} to improve probability calibration. A scalar temperature $T$ is found by grid search over $T \in [0.5, 3.0]$ with step 0.05, maximizing validation macro F1:

\begin{equation}
    \hat{p}(y|\mathbf{x}) = \text{softmax}\left(\frac{\hat{\mathbf{y}}}{T}\right)
    \label{eq:temp}
\end{equation}

Temperature scaling preserves the model's ranking of classes while adjusting confidence levels, making predictions more reliable for downstream decision-making.


% ═══════════════════════════════════════════════════════════════
% IV. EXPERIMENTAL SETUP
% ═══════════════════════════════════════════════════════════════
\section{Experimental Setup}

\subsection{Dataset}

We use the \textit{Sentiment Analysis for Mental Health} dataset from Kaggle~\cite{sarkar2023dataset}, which comprises 53,043 text samples across seven mental health categories. Table~\ref{tab:dataset} presents the class distribution.

\begin{table}[t]
    \centering
    \caption{Dataset Class Distribution}
    \label{tab:dataset}
    \begin{tabular}{lrr}
        \toprule
        \textbf{Class} & \textbf{Samples} & \textbf{Percentage} \\
        \midrule
        Normal               & 16,351 & 30.83\% \\
        Depression           & 15,404 & 29.04\% \\
        Suicidal             & 10,653 & 20.09\% \\
        Anxiety              &  3,888 &  7.33\% \\
        Bipolar              &  2,877 &  5.43\% \\
        Stress               &  2,669 &  5.03\% \\
        Personality Disorder &  1,201 &  2.26\% \\
        \midrule
        \textbf{Total}       & \textbf{53,043} & \textbf{100\%} \\
        \bottomrule
    \end{tabular}
\end{table}

The dataset exhibits significant class imbalance, with a 13.6:1 ratio between the largest (Normal) and smallest (Personality Disorder) classes. Texts are sourced from social media posts and online mental health forums.

\subsection{Data Augmentation}

To enrich the training data, we apply back-translation augmentation, generating 78,943 augmented samples for the three most critical classes: Depression (44,774), Suicidal (31,500), and Personality Disorder (2,669). Augmented data is used \textbf{only for training}; validation and test sets contain exclusively original samples to prevent near-duplicate leakage.

\subsection{Data Splitting}

The original dataset is split with stratified sampling:
\begin{itemize}[noitemsep]
    \item \textbf{Training set}: 80\% of original data + all augmented data
    \item \textbf{Validation set}: 10\% of original data
    \item \textbf{Test set}: 10\% of original data
\end{itemize}

A fixed random seed (42) ensures reproducibility.

\subsection{Implementation Details}

The model is implemented in PyTorch with the Hugging Face Transformers library. Training was conducted on a GPU-equipped system. Total training time was approximately 379 minutes (6.3 hours) for 10 epochs. Key hyperparameters are summarized in Table~\ref{tab:hyperparams}.

\begin{table}[t]
    \centering
    \caption{Hyperparameter Configuration}
    \label{tab:hyperparams}
    \begin{tabular}{lr}
        \toprule
        \textbf{Hyperparameter} & \textbf{Value} \\
        \midrule
        Pre-trained model       & RoBERTa-large \\
        Max sequence length     & 256 \\
        Batch size              & 16 \\
        Gradient accumulation   & 4 (eff. 64) \\
        Encoder LR ($\eta_{\text{bert}}$) & $2 \times 10^{-5}$ \\
        Head LR ($\eta_{\text{head}}$)    & $1 \times 10^{-4}$ \\
        LR decay ($\gamma$)     & 0.90 \\
        Weight decay ($\lambda$)& 0.02 \\
        Dropout                 & 0.30 / 0.15 \\
        Focal $\gamma$          & 2.0 \\
        Label smoothing $\epsilon$ & 0.05 \\
        Trainable layers        & 10 / 24 \\
        Warmup ratio            & 6\% \\
        Max epochs              & 20 \\
        Early stopping patience & 5 \\
        \bottomrule
    \end{tabular}
\end{table}


% ═══════════════════════════════════════════════════════════════
% V. RESULTS AND ANALYSIS
% ═══════════════════════════════════════════════════════════════
\section{Results and Analysis}

\subsection{Training Dynamics}

Table~\ref{tab:training} presents the epoch-by-epoch training progression. The model converges effectively over 10 epochs, with validation accuracy improving from 81.63\% (epoch 1) to 91.98\% (epoch 10).

\begin{table}[t]
    \centering
    \caption{Training History per Epoch}
    \label{tab:training}
    \resizebox{\columnwidth}{!}{
    \begin{tabular}{ccccccc}
        \toprule
        \textbf{Epoch} & \textbf{Tr Loss} & \textbf{Tr Acc} & \textbf{Val Acc} & \textbf{Tr F1} & \textbf{Val F1} & \textbf{Gap} \\
        \midrule
        1  & 0.2250 & 67.63\% & 81.63\% & 61.69\% & 80.32\% & $-$14.0\% \\
        2  & 0.1048 & 77.67\% & 84.70\% & 79.70\% & 84.21\% & $-$7.0\% \\
        3  & 0.0778 & 80.58\% & 85.38\% & 85.05\% & 84.72\% & $-$4.8\% \\
        4  & 0.0609 & 82.83\% & 87.06\% & 88.58\% & 87.36\% & $-$4.2\% \\
        5  & 0.0486 & 84.93\% & 87.99\% & 91.41\% & 88.26\% & $-$3.1\% \\
        6  & 0.0396 & 87.10\% & 89.75\% & 93.42\% & 89.35\% & $-$2.7\% \\
        7  & 0.0321 & 89.48\% & 90.26\% & 95.08\% & 89.45\% & $-$0.8\% \\
        8  & 0.0261 & 91.58\% & 90.99\% & 96.27\% & 90.26\% & +0.6\% \\
        9  & 0.0221 & 92.92\% & 91.67\% & 97.01\% & 90.62\% & +1.3\% \\
        10 & 0.0205 & 93.60\% & 91.98\% & 97.32\% & 90.73\% & +1.6\% \\
        \bottomrule
    \end{tabular}
    }
\end{table}

Several observations are noteworthy:

\begin{itemize}[noitemsep]
    \item \textbf{Underfitting to convergence}: In early epochs, the train-validation gap is negative (validation outperforms training), likely due to dropout regularization and augmented data noise in training. This reverses after epoch 7 as the model begins to memorize training examples.
    \item \textbf{Controlled overfitting}: The final train-val accuracy gap is only 1.6\%, well below the 8\% threshold commonly used to indicate problematic overfitting. This validates the effectiveness of partial fine-tuning and dropout.
    \item \textbf{Steady F1 improvement}: Macro F1 improves consistently, indicating balanced learning across all classes despite the imbalanced dataset.
\end{itemize}

\subsection{Test Set Performance}

On the held-out test set, the model achieves:

\begin{table}[H]
    \centering
    \caption{Final Test Set Results}
    \label{tab:test}
    \begin{tabular}{lr}
        \toprule
        \textbf{Metric} & \textbf{Score} \\
        \midrule
        Accuracy        & 91.67\% \\
        Macro F1        & 90.74\% \\
        Weighted F1     & 91.71\% \\
        Temperature $T$ & 1.00 \\
        \bottomrule
    \end{tabular}
\end{table}

The close agreement between macro F1 (90.74\%) and weighted F1 (91.71\%) indicates consistent performance across both majority and minority classes. The optimal temperature was found to be $T = 1.0$, indicating that the model's probability calibration is already well-tuned—a likely result of the combined effect of focal loss and label smoothing.

\subsection{Per-class Analysis}

The triple pooling approach and focal loss contribute to strong per-class performance. Majority classes (Normal, Depression, Suicidal) benefit from ample training data, while minority classes (Bipolar, Stress, Personality Disorder) benefit from the focal loss mechanism and class weighting, which prevents the model from learning to predict only majority classes.

The confusion patterns reveal expected clinical overlaps:
\begin{itemize}[noitemsep]
    \item Depression $\leftrightarrow$ Suicidal: These conditions share linguistic markers of hopelessness and despair.
    \item Anxiety $\leftrightarrow$ Stress: Worry-related language appears in both categories.
    \item Bipolar $\leftrightarrow$ Depression: The depressive phase of bipolar disorder produces language similar to clinical depression.
\end{itemize}

\subsection{Parameter Efficiency}

Table~\ref{tab:params} presents the parameter distribution of the model. By freezing 62\% of parameters, we reduce computational costs and regularize effectively.

\begin{table}[t]
    \centering
    \caption{Parameter Distribution}
    \label{tab:params}
    \begin{tabular}{lrr}
        \toprule
        \textbf{Component} & \textbf{Parameters} & \textbf{Trainable} \\
        \midrule
        Embedding layer            & $\sim$2M    & No \\
        Frozen encoder (0--13)     & $\sim$186M  & No \\
        Trainable encoder (14--23) & $\sim$133M  & Yes \\
        Classification head        & $\sim$1.6M  & Yes \\
        \midrule
        \textbf{Total}             & $\sim$\textbf{355M} & \textbf{134.6M (38\%)} \\
        \bottomrule
    \end{tabular}
\end{table}


% ═══════════════════════════════════════════════════════════════
% VI. DEPLOYED APPLICATION
% ═══════════════════════════════════════════════════════════════
\section{Deployed Application}

To facilitate practical use and demonstration, we developed an interactive web application using Streamlit. The application provides five functional modules:

\begin{enumerate}[noitemsep]
    \item \textbf{Single Text Classification}: Users input text and receive predicted class, confidence score, and a probability distribution chart across all seven classes. Crisis resources are automatically displayed when suicidal content is detected.
    \item \textbf{Batch Analysis}: Supports CSV upload or multi-line input for analyzing multiple texts simultaneously, with aggregate statistics and downloadable results.
    \item \textbf{Training Insights}: Interactive visualization of training curves (accuracy, F1, loss, overfitting gap) using Altair charts.
    \item \textbf{Architecture Explorer}: Visual representation of the model pipeline with detailed explanations of design decisions.
    \item \textbf{About}: Comprehensive project documentation including methodology, dataset information, and ethical considerations.
\end{enumerate}

The application supports both dark and light themes, runs entirely locally (no data is transmitted externally), and performs inference on CPU for deployment accessibility.


% ═══════════════════════════════════════════════════════════════
% VII. LIMITATIONS AND ETHICAL CONSIDERATIONS
% ═══════════════════════════════════════════════════════════════
\section{Limitations and Ethical Considerations}

\subsection{Limitations}

\begin{itemize}[noitemsep]
    \item \textbf{Dataset bias}: The training data is sourced primarily from English-language social media, which may not generalize to clinical settings, other languages, or cultural contexts.
    \item \textbf{Label overlap}: Comorbidities (e.g., depression with anxiety) are common in mental health but the dataset assumes single-label classification.
    \item \textbf{Temporal context}: The model analyzes individual text snippets without access to conversation history or longitudinal patterns.
    \item \textbf{Model size}: At 355M parameters, the model requires significant computational resources, potentially limiting deployment on edge devices.
    \item \textbf{Minority class performance}: Despite focal loss and augmentation, the rarest class (Personality Disorder, 2.26\% of data) may still exhibit lower reliability.
\end{itemize}

\subsection{Ethical Considerations}

This system is designed as a \textbf{research and educational tool only}. It is explicitly \textbf{not} intended to replace professional mental health assessment. Key ethical safeguards include:

\begin{itemize}[noitemsep]
    \item Clear disclaimers throughout the application interface
    \item Automatic provision of crisis hotline resources when suicidal ideation is detected
    \item Local-only processing — no user data is stored or transmitted
    \item Transparent reporting of model confidence to avoid false certainty
\end{itemize}

Mental health classification models carry inherent risks of misclassification. False negatives (failing to detect a condition) could lead to delayed intervention, while false positives may cause unnecessary alarm. Deployment in any clinical or screening context would require extensive validation, regulatory approval, and integration with human oversight.


% ═══════════════════════════════════════════════════════════════
% VIII. CONCLUSION AND FUTURE WORK
% ═══════════════════════════════════════════════════════════════
\section{Conclusion and Future Work}

This paper presented a comprehensive pipeline for multi-class mental health text classification using a fine-tuned RoBERTa-large model. The combination of partial fine-tuning, triple pooling, focal loss with label smoothing, and temperature scaling achieves a test accuracy of 91.67\% and macro F1-score of 90.74\% across seven mental health categories.

Key contributions include:
\begin{itemize}[noitemsep]
    \item A triple pooling mechanism that provides richer text representations than standard CLS-only approaches
    \item A training strategy combining focal loss, class weighting, and label smoothing that effectively handles 13.6:1 class imbalance
    \item Partial fine-tuning with layer-wise LR decay that reduces trainable parameters to 38\% of total while maintaining strong performance
    \item An interactive web application for real-time inference and model exploration
\end{itemize}

Future work may explore:
\begin{itemize}[noitemsep]
    \item \textbf{Multi-label classification} to capture comorbid conditions
    \item \textbf{Explainability} via attention visualization and SHAP values to identify which text segments drive predictions
    \item \textbf{Multilingual extension} by fine-tuning XLM-RoBERTa for cross-lingual mental health classification
    \item \textbf{Lightweight deployment} through knowledge distillation to a smaller model for mobile/edge inference
    \item \textbf{Longitudinal analysis} by incorporating temporal context from conversation histories
    \item \textbf{Domain-specific pre-training} using MentalBERT-style continued pre-training on mental health corpora
\end{itemize}


% ═══════════════════════════════════════════════════════════════
% REFERENCES
% ═══════════════════════════════════════════════════════════════
\begin{thebibliography}{99}

\bibitem{who2022}
World Health Organization, ``Mental disorders,'' \textit{WHO Fact Sheets}, June 2022. [Online]. Available: \url{https://www.who.int/news-room/fact-sheets/detail/mental-disorders}

\bibitem{calvo2017natural}
R.~A. Calvo, D.~N. Milne, M.~S. Hussain, and H. Christensen, ``Natural language processing in mental health applications using non-clinical texts,'' \textit{Natural Language Engineering}, vol.~23, no.~5, pp.~649--685, 2017.

\bibitem{coppersmith2015clpsych}
G.~Coppersmith, M.~Dredze, C.~Harman, and K.~Hollingshead, ``From ADHD to SAD: Analyzing the language of mental health on Twitter through self-reported diagnoses,'' in \textit{Proc. 2nd Workshop on Computational Linguistics and Clinical Psychology (CLPsych)}, 2015, pp.~1--10.

\bibitem{dechoudhury2013predicting}
M.~De~Choudhury, M.~Gamon, S.~Counts, and E.~Horvitz, ``Predicting depression via social media,'' in \textit{Proc. 7th International AAAI Conference on Weblogs and Social Media (ICWSM)}, 2013, pp.~128--137.

\bibitem{zirikly2019clpsych}
A.~Zirikly, P.~Resnik, O.~Uzuner, and K.~Hollingshead, ``CLPsych 2019 shared task: Predicting the degree of suicide risk in Reddit posts,'' in \textit{Proc. 6th Workshop on Computational Linguistics and Clinical Psychology (CLPsych)}, 2019, pp.~24--33.

\bibitem{matero2019suicide}
M.~Matero, A.~Idnani, Y.~Son, S.~Giorgi, H.~A. Schwartz, and S.~P. Ungar, ``Suicide risk assessment with multi-level dual-context language and BERT,'' in \textit{Proc. 6th Workshop on Computational Linguistics and Clinical Psychology (CLPsych)}, 2019, pp.~39--44.

\bibitem{devlin2019bert}
J.~Devlin, M.~W. Chang, K.~Lee, and K.~Toutanova, ``BERT: Pre-training of deep bidirectional transformers for language understanding,'' in \textit{Proc. NAACL-HLT}, 2019, pp.~4171--4186.

\bibitem{liu2019roberta}
Y.~Liu \textit{et al.}, ``RoBERTa: A robustly optimized BERT pretraining approach,'' \textit{arXiv preprint arXiv:1907.11692}, 2019.

\bibitem{ji2022mentalbert}
S.~Ji, T.~Zhang, L.~Ansari, J.~Fu, P.~Tiwari, and E.~Cambria, ``MentalBERT: Publicly available pretrained language models for mental healthcare,'' in \textit{Proc. LREC}, 2022, pp.~7184--7190.

\bibitem{lin2017focal}
T.-Y. Lin, P.~Goyal, R.~Girshick, K.~He, and P.~Doll\'{a}r, ``Focal loss for dense object detection,'' in \textit{Proc. IEEE ICCV}, 2017, pp.~2980--2988.

\bibitem{cui2019class}
Y.~Cui, M.~Jia, T.-Y. Lin, Y.~Song, and S.~Belongie, ``Class-balanced loss based on effective number of samples,'' in \textit{Proc. IEEE CVPR}, 2019, pp.~9268--9277.

\bibitem{reimers2019sentencebert}
N.~Reimers and I.~Gurevych, ``Sentence-BERT: Sentence embeddings using Siamese BERT-networks,'' in \textit{Proc. EMNLP-IJCNLP}, 2019, pp.~3982--3992.

\bibitem{conneau2017supervised}
A.~Conneau, D.~Kiela, H.~Schwenk, L.~Barrault, and A.~Bordes, ``Supervised learning of universal sentence representations from natural language inference data,'' in \textit{Proc. EMNLP}, 2017, pp.~670--680.

\bibitem{guo2017calibration}
C.~Guo, G.~Pleiss, Y.~Sun, and K.~Q. Weinberger, ``On calibration of modern neural networks,'' in \textit{Proc. ICML}, 2017, pp.~1321--1330.

\bibitem{sarkar2023dataset}
S.~Sarkar, ``Sentiment analysis for mental health,'' \textit{Kaggle}, 2023. [Online]. Available: \url{https://www.kaggle.com/datasets/suchintikasarkar/sentiment-analysis-for-mental-health}

\end{thebibliography}

\end{document}
